\chapter{Comparing Operators Responsibly}
\index{operator comparison|textbf}

\label{ch:mathematical-foundations}

\chapterstatus{expository mathematical framework; no numerical or scientific
result is claimed.}

This chapter fixes the distinctions needed to discuss projection,
representation, alignment, subtraction, and reduction. It does not prescribe a
particular alignment algorithm or reduced model class. Those choices belong to
the applicable methodological specifications.

\section{State spaces and operators}
\index{state space}
\index{operator!abstract}


Let $\mathcal H$ denote the parent one-particle state space and let
$\widehat H$ denote the selected Kohn--Sham operator on its stated domain in
$\mathcal H$. The physical model, the operator $\widehat H$, and any finite
matrix used by software are different objects. The notation $H$ is reserved
below for a finite matrix only after its represented space and basis have been
identified.

A finite retained subspace $\mathcal S\subset\mathcal H$ has dimension $N$ and
orthogonal projector
\begin{equation}
  P_{\mathcal S}:\mathcal H\longrightarrow\mathcal S.
\end{equation}
When the required domain conditions hold, the compressed operator on the
retained space is
\begin{equation}
  \widehat H_{\mathcal S}
  =
  \left.P_{\mathcal S}\widehat H P_{\mathcal S}\right|_{\mathcal S}
  :\mathcal S\longrightarrow\mathcal S.
\end{equation}
Selecting $\mathcal S$ is a model or representation decision. Writing the
coordinates of $\widehat H_{\mathcal S}$ in a basis is a separate operation.
Neither operation is synonymous with Wannier localization or with truncation of
matrix entries.

\section{Finite matrix representations}
\index{matrix representation}
\index{represented operator}


Let
\begin{equation}
  \mathcal B=\{\lvert b_i\rangle\}_{i=0}^{N-1}
\end{equation}
be an ordered orthonormal basis for $\mathcal S$. The represented matrix is
\begin{equation}
  H^{(\mathcal B)}_{ij}
  =
  \langle b_i\lvert\widehat H_{\mathcal S}\rvert b_j\rangle,
  \qquad
  H^{(\mathcal B)}\in\mathbb C^{N\times N}.
\end{equation}
The array of complex numbers is not a complete represented operator by itself.
Its interpretation can depend on

\begin{itemize}
  \item the identity and dimension of $\mathcal S$;
  \item basis identity, orbital convention, and ordering;
  \item geometry, lattice vectors, site labels, and boundary convention;
  \item units and normalization;
  \item scalar energy reference;
  \item spin representation; and
  \item gauge, phase, or other coordinate conventions.
\end{itemize}

Equal shape and data type do not establish equality of represented meaning.
Likewise, common provenance identifies a source calculation but does not prove
that two records use compatible coordinates.

\section{Basis transformations and gauge}
\index{basis transformation}
\index{gauge}


Suppose a second ordered orthonormal basis is related to the first by
$U\in U(N)$. With the convention
\begin{equation}
  \lvert b'_j\rangle
  =
  \sum_{i=0}^{N-1}\lvert b_i\rangle U_{ij},
\end{equation}
the matrix coordinates transform as
\begin{equation}
  H^{(\mathcal B')}
  =
  U^{\dagger}H^{(\mathcal B)}U.
\end{equation}
This transformation changes coordinates without changing the represented
operator. Individual matrix entries and entrywise differences are therefore
basis dependent.

For an isolated retained Bloch subspace, momentum-dependent unitary freedom
similarly changes the Wannier basis generated from that subspace. Localization
selects or optimizes a gauge; it does not remove the need to record the selected
frame. Equal band energies do not determine Wannier functions or provide a
physically distinguished alignment between two calculations.

\citationtodo{medium priority}{Check the specialization from finite-dimensional
basis changes to Bloch/Wannier gauge and localization against Marzari et al.
(2012) and Souza--Marzari--Vanderbilt (2001), bibliography keys
\texttt{marzari2012} and \texttt{souza2001} (already present).}

\section{Spectra and operator information}
\index{spectrum!operator information}


The spectrum of a finite Hermitian matrix is invariant under unitary
conjugation. Spectral comparison is consequently valuable but incomplete for
operator comparison. Even when two Hermitian matrices are isospectral, the
spectrum alone does not identify which basis vectors, sites, orbitals, spin
channels, or spatial regions correspond. It supplies no distinguished physical
alignment map.

Conversely, a reduced model may approximate selected eigenvalues while
retaining different onsite terms or couplings in the coordinates used to
represent a later perturbation. A spectral acceptance rule and an aligned
operator acceptance rule are therefore separate constraints. Their relationship
must be investigated rather than assumed.

\section{Alignment between retained spaces}
\index{alignment!retained spaces}
\index{identification map}


Let $\mathcal S_{\mathrm r}$ and $\mathcal S_{\mathrm c}$ denote reference and
candidate retained spaces of equal finite dimension. An alignment is an
explicitly defined isometric identification
\begin{equation}
  A:\mathcal S_{\mathrm r}\longrightarrow\mathcal S_{\mathrm c},
  \qquad
  A^{\dagger}A=I_{\mathcal S_{\mathrm r}}.
\end{equation}
When the spaces have equal dimension and $A$ is onto, $A$ is unitary. The
candidate operator pulled back to the reference space is
\begin{equation}
  \widehat H_{\mathrm c\rightarrow r}
  =
  A^{\dagger}\widehat H_{\mathrm c}A
  :\mathcal S_{\mathrm r}\longrightarrow\mathcal S_{\mathrm r}.
\end{equation}
The direction of $A$, its coordinate convention, and the information used to
construct it must be stated. Orbital labels, matrix dimension, or similar
spectra do not substitute for this map.

If no justified alignment exists, coordinate-wise subtraction is undefined for
the intended physical comparison. One may instead use a declared invariant
comparison, but the invariant and the claim it supports must be named
explicitly.

\section{Energy references}
\index{energy reference}


An energy-valued operator is unchanged in its eigenvectors but shifted in its
spectrum when a scalar multiple of the identity is added. For a recorded energy
zero $E_0$, define
\begin{equation}
  \overline H=H-E_0I.
\end{equation}
Reference and candidate operators must use a common physical energy convention
before their difference is interpreted. A valence-band maximum used to plot
bulk bands and an electrostatic estimator used to align pristine and doped
supercells are not automatically the same convention.

The physical specification requires a common scalar energy zero before
bulk--dopant subtraction, while the operational estimator and its uncertainty
belong to a later alignment stage. Until that estimator is supplied, the
requirement is fixed but the corresponding impurity difference is not yet a
calculated project result.

\citationtodo{high priority}{Check the distinction between plotting references
and electrostatic defect-supercell alignment against Freysoldt et al. (2009,
2014), bibliography keys \texttt{freysoldt2009} and \texttt{freysoldt2014}
(already present).}

\section{Represented differences}
\index{operator difference!represented}


After compatibility and alignment have been established, define the signed
represented difference by an explicit operand order. For example,
\begin{equation}
  \Delta H
  =
  \overline H_{\mathrm c\rightarrow r}
  -
  \overline H_{\mathrm r}.
\end{equation}
Reversing the operands reverses the sign. A difference matrix is not
intrinsically an impurity operator, perturbation, approximation error, or
observable. That interpretation follows only from the identities of the two
operators, their common-space construction, and the declared scientific use.

For pristine and doped systems, compatibility includes at least the represented
state spaces, matrix dimension, basis ordering, geometry and site mapping,
units, energy reference, spin convention, and gauge alignment. An unresolved
item is a stop condition for direct matrix subtraction.

\section{Hermiticity and residuals}
\index{Hermiticity}
\index{operator residual}


A represented Hamiltonian is expected to be Hermitian within a stated numerical
criterion. One possible absolute residual is
\begin{equation}
  \varepsilon_{\mathrm H}
  =
  \max_{i,j}\left|H_{ij}-H_{ji}^{*}\right|.
\end{equation}
The residual, its units, any normalization scale, and its acceptance tolerance
are distinct quantities. Passing a Hermiticity tolerance verifies only that
declared numerical criterion; it does not validate the physical model or the
metadata used to interpret the matrix.

For a compatible difference $\Delta H$, examples of absolute residual measures
include
\begin{align}
  \varepsilon_{\max} &= \max_{i,j}|\Delta H_{ij}|,\\
  \varepsilon_{\mathrm F} &= \|\Delta H\|_{\mathrm F},\\
  \varepsilon_{2} &= \|\Delta H\|_{2}.
\end{align}
A normalized residual must identify its reference scale and define behavior
when that scale vanishes. Numerical implementations must also handle nonfinite
values and overflow explicitly rather than silently reporting an infinite or
undefined metric.

\section{Approximation and evidence boundaries}
\index{approximation!evidence boundary}


Projection error, discretization error, fitting error, truncation error, and
parent-model error arise from different operations. This monograph groups them
only when an owning mathematical contract defines a compatible relationship.
In particular:

\begin{itemize}
  \item construction and invariant checks provide software verification;
  \item agreement with independently derived matrix identities or norms may
    provide numerical verification;
  \item comparison with independent physical evidence may provide scientific
    validation for a declared use; and
  \item sensitivity analysis or propagation of declared uncertainties belongs
    to uncertainty quantification.
\end{itemize}

No one category implies another. These distinctions determine how later
chapters may describe projected Hamiltonians, fitted model classes, and
impurity-operator residuals without turning coordinate agreement into an
unsupported physical conclusion. Chapter~\ref{ch:proof-program} records the
corresponding proof obligations and their current analytical and mechanization
status.
